problem:  
Let \(\mathcal{C}_{\text{HBC}}\), \(\mathcal{C}_{\text{OBC}}\), and \(\mathcal{C}_{\text{Ric}^\perp}\) denote the standard unitary-invariant convex cones in the space of Kähler curvature tensors on a Hermitian vector space \(V\), corresponding respectively to nonnegativity of **holomorphic bisectional curvature**, **orthogonal bisectional curvature**, and **orthogonal Ricci curvature**.  

Define a convex cone \(\mathcal{C}\) of Kähler curvature tensors to be **Frankel-effective** if, for every compact Kähler manifold \((M^{2n}, g, J)\) with \(\operatorname{Rm}(M)\in \mathcal{C}\) pointwise, the following holds:

> Every pair of closed complex submanifolds \(X^p, Y^q \subset M\) with \(p+q \ge n\) satisfies \(X\cap Y \neq \varnothing.\)

**Open problem:**  
Classify all **maximal unitary-invariant convex cones** \(\mathcal{C}_{\mathrm{Frankel}}\) in the space of Kähler curvature tensors that are Frankel-effective, i.e., that imply the intersection property above but are not properly contained in any other unitary-invariant convex cone with the same property. Specifically:

1. Does there exist a maximal \(\mathcal{C}_{\mathrm{Frankel}}\) strictly larger than the bisectional-positivity cone \(\mathcal{C}_{\text{HBC}}\)?  
2. Do the known curvature positivity cones—such as nonnegative quadratic bisectional curvature (\(\mathcal{C}_{\text{QBC}}\)) or positive orthogonal Ricci curvature (\(\mathcal{C}_{\text{Ric}^\perp}\))—lie within any Frankel-effective maximal cone?  
3. Can one characterize \(\mathcal{C}_{\mathrm{Frankel}}\) analytically in terms of the **nonnegativity of the index form** arising from the second variation of distance between two complex submanifolds?  

Why is it a "good" problem:  
This formulation removes the ad hoc curvature mixtures and directly addresses the **structural geometric core** of Frankel-type phenomena. It poses a rigorous classification problem: to identify all curvature positivity cones guaranteeing the submanifold-intersection property. Such a problem bridges several high-level areas—representation theory of \(U(n)\) acting on curvature tensors, convex geometry of curvature cones, and analytic methods in the second variation of geodesic distance.

The question is simple to state yet conceptually profound. It seeks to delineate the precise **interface between local curvature positivity and global topological rigidity**. A resolution would unify known partial results (Frankel’s theorem, QBC, Ric\(^\perp\)-positivity studies) and chart the full boundary in curvature space separating those Kähler geometries that enforce intersection from those that may admit separation. This interplay of curvature cone geometry and global intersection theory exemplifies an elegant and deep research direction—precisely the hallmark of a “good” mathematical problem.